\documentclass[12pt]{article}
\usepackage{graphicx}
\usepackage{setspace}
\usepackage{indentfirst}
%\usepackage[margin=1in]{geometry}
\usepackage[letterpaper,top=2cm,bottom=2cm,left=3cm,right=3cm,marginparwidth=1.75cm]{geometry}

% Useful packages
\usepackage{amsmath}
\usepackage{amssymb}
\usepackage{graphicx}
\usepackage{booktabs}
\usepackage{longtable}
\usepackage[colorlinks=true, allcolors=blue]{hyperref}

\title{BIOSTAT M235 -- Final Project\\
\large Does Health Insurance Affect Whether U.S.\ Adults Receive a
Blood-Sugar Screening? A Causal Analysis of NHANES 2017--2020}
\author{Amanda Millatt \and Janvi Bharucha \and Sujit Silas Armstrong}
\date{\today}

\begin{document}
\doublespacing

\maketitle

\newpage

\section{Explanation}

\subsection{Introduction} \label{subsec:intro}

More than 38 million Americans have diabetes, with an estimated 90--95\% of
cases being type~2 diabetes \cite{cdc2024ndsr}. In addition, over one third of
American adults have prediabetes, and many of those will progress to type~2
diabetes within three to five years \cite{cdc2024ndsr}. Because early detection
depends on a blood-sugar screening test, understanding what determines whether
an individual is screened is an important public-health question.

Health insurance in the United States has been established as a critical
determinant of health, including diabetes-related outcomes. Insurance coverage
facilitates early diagnosis through access to screening, whereas uninsured
individuals may face greater difficulty accessing care, leading to delays in
diagnosis and reduced treatment adherence \cite{bolarinwa2025}. Prior NHANES
analyses reinforce this: insured patients are substantially more likely to meet
diabetes care-quality benchmarks, with one study reporting that insured patients
had 2.73 times the odds of controlled blood pressure relative to uninsured
patients \cite{casagrande2016}. A further NHANES analysis found that having
health insurance was associated with significantly lower odds of undiagnosed
prediabetes and type~2 diabetes \cite{mahoney2020}. The CDC has likewise
identified uninsured people as being at elevated risk of undetected diabetes,
noting that more than one in four adults with diabetes do not know they have it,
and that eight in ten Americans with prediabetes are unaware of their risk
\cite{cdc2023screening}.

While Mahoney et al.\ \cite{mahoney2020} demonstrated that insurance is
associated with lower odds of \emph{undiagnosed} diabetes, their outcome was
undiagnosed disease status. Our analysis instead targets the step \emph{before}
diagnosis: whether an individual receives a blood-sugar screening test at all.
We frame this as a causal question and estimate the average treatment effect
(ATE) of insurance coverage on screening using three complementary
estimators: survey-weighted logistic regression, inverse-probability weighting
(IPW), and a doubly robust (AIPW) estimator. We then run a battery of
sensitivity analyses.

\subsection{Set Up} \label{subsec:setup}

\paragraph{Data source.} We use the National Health and Nutrition Examination
Survey (NHANES) 2017--March~2020 pre-pandemic cycle. The exposure is the
self-reported binary indicator \texttt{HIQ011} (``Are you covered by health
insurance or some other kind of health-care plan?''), coded as insured versus
uninsured. The outcome is \texttt{DIQ180}, self-reported receipt of a blood-sugar
screening test in the past three years. Because NHANES skips \texttt{DIQ180} for
respondents already diagnosed with diabetes, the analytic population consists of
\emph{undiagnosed} adults by construction.

\paragraph{Causal estimand.} Let $A$ denote insurance coverage ($1$ = insured,
$0$ = uninsured) and $Y$ the screening outcome. We target the average treatment
effect on the risk-difference scale,
\begin{equation}
\tau \;=\; \mathbb{E}\!\left[Y^{1}\right] - \mathbb{E}\!\left[Y^{0}\right],
\label{eq:ate}
\end{equation}
the change in the population probability of screening that would result if every
adult were insured versus uninsured. Identification of $\tau$ from observational
data rests on three assumptions \cite{hernan2020}:
\begin{enumerate}
\item \textbf{Consistency:} the observed outcome equals the potential outcome
under the observed exposure.
\item \textbf{Conditional exchangeability (no unmeasured confounding):}
$Y^{a} \perp A \mid L$ for the measured covariate set $L$.
\item \textbf{Positivity:} $0 < \Pr(A = a \mid L = l) < 1$ for every covariate
stratum $l$ in the population.
\end{enumerate}

\paragraph{Confounders.} The covariate set $L$ was selected using a directed
acyclic graph (DAG) constructed from prior literature (see the project
repository, \texttt{nhanes\_dag.Rmd}), choosing variables that plausibly affect
both insurance coverage and the propensity to be screened. Table~\ref{tab:covs}
lists the covariates and their rationale. Notably, variables that lie on the
causal path from insurance to screening (e.g.\ having a usual source of care or
the number of health-care visits) were deliberately \emph{excluded} as mediators
rather than confounders.

{\singlespacing
\begin{longtable}{p{3cm} p{10.5cm}}
\caption{Covariates included in the adjustment set $L$.}
\label{tab:covs}\\
\toprule
\textbf{Covariate} & \textbf{Rationale}\\
\midrule
\endfirsthead
\multicolumn{2}{c}{\tablename\ \thetable\ (continued)}\\
\toprule
\textbf{Covariate} & \textbf{Rationale}\\
\midrule
\endhead
Age & Older adults are more likely to be insured (e.g.\ Medicare eligibility at
65) and are more frequently referred for screening due to higher baseline
diabetes risk.\\
Sex & Men and women may differ in healthcare-seeking behavior and
insurance-coverage patterns.\\
Race/ethnicity & Racial and ethnic minorities have systematically lower
insurance rates and differ in diabetes risk and screening likelihood.\\
Education & Higher education is associated with greater health literacy,
employer-sponsored insurance, and engagement with preventive care.\\
Marital status & Married/partnered individuals are more likely to be covered
under a spouse's plan and to have higher healthcare utilization.\\
Country of birth & Foreign-born individuals face barriers to obtaining insurance
and accessing preventive screening.\\
Interview language & Proxy for acculturation.\\
Income-to-poverty ratio & Lower income is strongly associated with both lack of
coverage and reduced preventive-care utilization.\\
BMI & Elevated BMI is a clinical risk factor for type~2 diabetes, so providers
are more likely to order screening regardless of insurance.\\
Self-rated health & Individuals rating their health poor/fair may seek care
more, but may also be less likely to be insured.\\
Hypertension & A major comorbidity co-occurring with insulin resistance and
metabolic syndrome that itself prompts screening.\\
Ever smoked $\geq$100 cigarettes & Associated with socioeconomic disadvantage,
lower insurance rates, and insulin resistance.\\
Ever had alcohol & Heavy alcohol use is linked to metabolic dysregulation,
affecting glucose metabolism and screening decisions.\\
\bottomrule
\end{longtable}
}

\paragraph{Sample and weighting.} The analytic sample was restricted to adults
aged $\geq 20$ years who answered both the insurance and screening questions.
Respondents who completed the interview but not the physical exam (zero MEC
weight) were excluded because BMI is measured in the exam. All estimates
incorporate the NHANES complex sampling design, namely primary sampling units
(\texttt{SDMVPSU}), strata (\texttt{SDMVSTRA}), and MEC exam weights
(\texttt{WTMECPRP}), via the \texttt{survey} package \cite{lumley2004}, so that
estimates are nationally representative. Covariate missingness was handled by
complete-case analysis in the primary models and probed by multiple imputation
in the sensitivity analyses. The full exclusion cascade is summarized in the
STROBE diagram in the repository.

\subsection{Methodology and Main Analysis} \label{subsec:meth_analysis}

We estimated $\tau$ three ways. Triangulating across estimators with different
modeling assumptions guards against any single method's misspecification.

\subsubsection{Inverse Probability Weighting (IPW)}

Inverse-probability-of-treatment weighting models the propensity score
$e(L) = \Pr(A = 1 \mid L)$ by logistic regression and reweights individuals by
the inverse of their probability of the observed exposure, creating a
pseudo-population in which exposure is independent of the measured covariates
\cite{robins2000}. Stabilized weights were used to reduce variance and were
combined with the NHANES survey weights (\texttt{WTMECPRP}) for national
representativeness. For each unit $i$, the stabilized weight is
\begin{equation}
sw_i = \frac{A_i\,\Pr(A = 1)}{e(L_i)} + \frac{(1 - A_i)\,\Pr(A = 0)}{1 - e(L_i)},
\end{equation}
where $\Pr(A = 1)$ is the marginal probability of exposure. The propensity-score
model was a binomial logistic regression over the covariates in
Table~\ref{tab:covs}, fit with the NHANES sample-design weights. To avoid
instability and convergence issues, the weights were rescaled, and observations
with missing covariates were dropped. Covariate balance before and after
weighting was assessed by plotting absolute standardized mean differences in a
Love plot, and a density plot of the estimated propensity scores by insurance
status was used to evaluate overlap and verify the positivity assumption.

The IPW-adjusted ATE ($\hat{\tau}_{\mathrm{IPW}}$) was computed as the difference
between the weighted potential-outcome means. Because the analytic standard
errors do not account for the estimation uncertainty of the propensity score,
non-parametric bootstrapping was conducted with $R = 500$ replications: the
baseline dataset was resampled with replacement, and each replication
re-estimated the propensity score, recomputed the stabilized IPW weights,
multiplied them by the survey-design weights, and rescaled the final weights to
mean one before fitting the outcome model.

\subsubsection{Regression Analysis}

A survey-weighted logistic regression model was used as the initial method,
since the outcome is binary. Three models were fit with increasing adjustment.
\emph{Model~1} was unadjusted. \emph{Model~2} added sociodemographic covariates
(age, sex, race/ethnicity, education, income-to-poverty ratio, marital status,
country of birth, interview language). \emph{Model~3} additionally adjusted for
health status and behavior (BMI, self-rated health, hypertension, smoking,
alcohol). The dispersion parameter was near one, supporting a binomial model.
Center weights (\texttt{WTMECPRP}) were used for national representativeness, and
observations with missing covariates were dropped. A Wald test assessed the
significance of the insurance term in each model, and the ATE was obtained by
g-computation (standardization). Results appear in Table~\ref{tab:regression}.

\subsubsection{Doubly Robust Estimation (AIPW)}

As the primary analysis we used the augmented inverse-probability-weighted
(AIPW) estimator, which is \emph{doubly robust}: it combines a propensity model
$e(L)$ with an outcome model $m_a(L) = \mathbb{E}[Y \mid A = a, L]$ and remains
consistent for $\tau$ if \emph{either} model (but not necessarily both) is
correctly specified \cite{robins1994, bang2005, funk2011}. For each subject $i$
we form the pseudo-outcomes
\begin{align}
\hat\mu_1^{(i)} &= m_1(L_i) + \frac{A_i}{e(L_i)}\bigl(Y_i - m_1(L_i)\bigr),\\
\hat\mu_0^{(i)} &= m_0(L_i) + \frac{1-A_i}{1-e(L_i)}\bigl(Y_i - m_0(L_i)\bigr),
\end{align}
whose survey-weighted means estimate $\mathbb{E}[Y^{1}]$ and $\mathbb{E}[Y^{0}]$;
the ATE is $\hat\tau = \overline{\hat\mu_1} - \overline{\hat\mu_0}$. Both nuisance
models were logistic regressions fit with the NHANES exam weights over the
covariates in Table~\ref{tab:covs}, excluding alcohol use; alcohol is omitted
from the primary doubly robust specification to limit complete-case loss, and the
sensitivity analysis (Table~\ref{tab:sens}) confirms that re-including it leaves
the estimate essentially unchanged ($0.189$ vs.\ $0.187$). (To avoid numerical instability, the exam
weights were rescaled to mean one within the nuisance fits; this leaves the
coefficients unchanged but prevents the spurious separation that very large case
weights induce.) Design-based confidence intervals for the risk difference, risk
ratio, and marginal odds ratio were obtained by the delta method applied to the
survey-weighted potential-outcome means, so that the NHANES clustering and
stratification are respected.

\subsubsection{Main Results}

After excluding interview-only respondents, the regression analytic sample was
6{,}894 adults: 5{,}694 (82.6\%) insured and 1{,}200 (17.4\%) uninsured. The
weighted screening rate was already markedly different by group: 57\% among
insured adults versus 34\% among uninsured adults. The doubly robust analysis,
which requires complete data on all twelve covariates, used 5{,}822 adults and
estimated $\mathbb{E}[Y^{1}] = 0.56$ and $\mathbb{E}[Y^{0}] = 0.37$.

\begin{table}[h]
\centering
\caption{Survey-weighted logistic regression: odds ratio (OR) for insurance and
g-computation ATE, by level of adjustment.}
\label{tab:regression}
\begin{tabular}{lcccc}
\toprule
\textbf{Model} & \textbf{OR} & \textbf{95\% CI} & \textbf{$p$-value} & \textbf{ATE}\\
\midrule
1: Unadjusted              & 2.52 & (2.11, 3.01) & $<0.001$ & 0.225\\
2: Demographically adjusted & 1.95 & (1.51, 2.52) & $0.0002$ & 0.156\\
3: Fully adjusted          & 1.94 & N/A          & $0.108$  & 0.148\\
\bottomrule
\end{tabular}
\end{table}

In the unadjusted model, insured adults had about 2.5 times the odds of being
screened (OR $= 2.52$, 95\% CI $2.11$--$3.01$, $p < 0.001$). Adjusting for
sociodemographic covariates attenuated this to OR $= 1.95$ (95\% CI
$1.51$--$2.52$, $p = 0.0002$); adding health and behavioral variables barely
changed it (OR $= 1.94$). Within Model~2, age (OR $= 1.02$ per year, $p =
0.0003$) and education (some college OR $= 1.41$; college graduate OR $= 1.45$)
were significant predictors of screening. Wald tests confirmed insurance as a
significant predictor in Model~1 ($F = 115.5$, $p < 0.001$) and Model~2 ($F =
35.0$, $p = 0.0002$). Because NHANES masks its true PSU identifiers, very few
residual degrees of freedom remain once Model~3's covariates are included,
rendering its confidence intervals unreliable; we therefore rely on Model~2 for
inference and report only Model~3's point estimates and $p$-values.

The three estimators agree closely (Table~\ref{tab:methods}). The standalone IPW
analysis yielded an adjusted ATE ($\hat{\tau}_{\mathrm{IPW}}$) of $0.191$,
corresponding to a marginal OR of $2.17$ (log-odds $0.777$, $p < 0.001$) within
the weighted pseudo-population, with a bootstrap 95\% CI for the risk difference
of $(0.108, 0.269)$ from $R = 500$ replications. The doubly robust risk
difference was $0.187$ (95\% CI $0.124$--$0.249$), corresponding to a marginal OR
of about $2.1$ and a risk ratio of $1.50$. All three methods indicate that
insurance raises the probability of blood-sugar screening by roughly 15--19
percentage points.

\begin{table}[h]
\centering
\caption{Average treatment effect (risk difference) across the three estimators.}
\label{tab:methods}
\begin{tabular}{lcc}
\toprule
\textbf{Method} & \textbf{ATE (risk difference)} & \textbf{95\% CI}\\
\midrule
Regression (g-computation, Model 2) & 0.156 & N/A\\
Inverse probability weighting       & 0.191 & (0.108, 0.269)\\
Doubly robust (AIPW)                & 0.187 & (0.124, 0.249)\\
\bottomrule
\end{tabular}
\end{table}

\section{Discussion} \label{sec:discussion}

\subsection{Advantages} \label{subsec:advantages}

Using three estimators with different assumptions strengthens our conclusions:
agreement across methods makes the finding robust to any one method's
misspecification. Logistic regression is straightforward to implement and
interpret, yields odds ratios comparable to prior literature, and uses the full
analytic sample. IPW explicitly targets the ATE and separates the design stage
(propensity modeling and balance checking) from outcome analysis. The doubly
robust estimator offers the strongest theoretical guarantee, requiring only one
of the two nuisance models to be correct, and, unlike regression alone, directly
estimates a marginal causal estimand on the risk-difference scale. All analyses
incorporate the NHANES survey design, so estimates generalize to the U.S.\ adult
population.

\subsection{Limitations} \label{subsec:limitations}

Several limitations temper the interpretation. First, both exposure and outcome
are self-reported and subject to recall and social-desirability bias. Second,
NHANES is cross-sectional, so insurance status and screening are measured
contemporaneously, and reverse pathways cannot be fully excluded. Third,
positivity is strained at the top of the propensity distribution: adults aged
65+ are nearly universally insured through Medicare ($\approx 98\%$), so there
are few comparable uninsured older adults, and effect estimates in that region
rely on model extrapolation. Fourth, complete-case analysis assumes data are
missing at random; income-to-poverty ratio and alcohol use carry the most
missingness. Finally, residual confounding from unmeasured factors (e.g.\
provider behavior, local health-system characteristics) remains possible, which
motivates the formal sensitivity analysis below.

\subsection{Comments} \label{subsec:finalcomments}

Across regression, IPW, and doubly robust estimation, having health insurance is
associated with a 15--19 percentage-point increase in the probability of
receiving a blood-sugar screening among undiagnosed U.S.\ adults. This extends
Mahoney et al.\ \cite{mahoney2020} upstream of diagnosis: insurance appears to
shape not only whether diabetes is detected, but whether the screening step that
enables detection occurs at all.

\section{Extension: Sensitivity Analysis} \label{sec:Extension}

\subsection{Motivation} \label{subsec:ExtMotivation}

The causal interpretation of $\tau$ depends on assumptions that cannot be fully
verified from the data: no unmeasured confounding, positivity, a correctly
specified adjustment set, and an appropriate treatment of missing data. As an
extension, we quantify how fragile the headline doubly robust estimate is to
plausible violations of each. All checks re-use the AIPW estimator as the
workhorse and are implemented in \texttt{sensitivity/sensitivity\_analysis.Rmd}.

\subsection{Method} \label{subsec:ExtMethod}

We examined four threats. (1) \emph{Unmeasured confounding} was quantified with
the E-value \cite{vanderweele2017}, the minimum strength of association (on the
risk-ratio scale) that an unmeasured confounder would need with both insurance
and screening to explain away the estimate. (2) \emph{Positivity} was probed by
re-estimating the ATE after truncating the propensity score to progressively
tighter bounds and after trimming the sample to the region of common support.
(3) \emph{Adjustment-set sensitivity} compared the primary twelve-covariate set
against adding alcohol use, a demographic/SES-only set, and a minimal set.
(4) \emph{Missing data} was addressed by multiple imputation by chained
equations (five imputations), re-estimating the AIPW within each completed
dataset and pooling by Rubin's rules \cite{rubin1987}.

\subsection{Results} \label{subsec:ExtResults}

The primary AIPW risk difference of $0.187$ corresponds to an E-value of
$\mathbf{2.37}$, with the lower confidence bound's E-value equal to $1.82$. That
is, an unmeasured confounder would need to be associated with both insurance and
screening by a risk ratio of at least $2.37$, beyond all measured covariates,
to fully explain away the effect, which is a substantial degree of confounding. Positivity restrictions (Table~\ref{tab:sens}) attenuated the
estimate only modestly, to $0.15$--$0.16$, with every confidence interval still
excluding zero. Varying the adjustment set produced risk differences from
$0.150$ (minimal) to $0.189$ (adding alcohol), all significant. Finally,
multiple imputation gave a pooled risk difference of $0.168$ (95\% CI
$0.113$--$0.223$), close to the complete-case value of $0.187$, indicating that
complete-case analysis did not materially bias the estimate.

\begin{table}[h]
\centering
\caption{Sensitivity of the doubly robust risk difference (RD).}
\label{tab:sens}
\begin{tabular}{llcc}
\toprule
\textbf{Check} & \textbf{Specification} & \textbf{RD} & \textbf{95\% CI}\\
\midrule
Positivity      & No trimming               & 0.187 & (0.124, 0.249)\\
                & Truncate PS $[0.05,0.95]$ & 0.161 & (0.124, 0.197)\\
                & Truncate PS $[0.10,0.90]$ & 0.160 & (0.131, 0.189)\\
                & Trim to PS $[0.05,0.95]$  & 0.150 & (0.087, 0.212)\\
\midrule
Adjustment set  & Primary (12 covariates)   & 0.187 & (0.124, 0.249)\\
                & $+$ alcohol use           & 0.189 & (0.124, 0.254)\\
                & Demographic/SES only      & 0.169 & (0.103, 0.235)\\
                & Minimal                   & 0.150 & (0.090, 0.210)\\
\midrule
Missing data    & Complete case             & 0.187 & (0.124, 0.249)\\
                & Multiple imputation ($m=5$) & 0.168 & (0.113, 0.223)\\
\bottomrule
\end{tabular}
\end{table}

\subsection{Discussion} \label{subsec:ExtDiscuss}

The headline finding is stable across all four checks: the effect direction and
rough magnitude persist under tighter positivity restrictions, alternative
adjustment sets, and multiple imputation, and the E-value indicates that only a
fairly strong unmeasured confounder could overturn it. The largest single
movement comes from restricting positivity, which both attenuates the estimate
and narrows it toward the region where insured and uninsured adults are
genuinely comparable, consistent with the Medicare-driven overlap concern noted
among older adults. We therefore interpret $0.15$--$0.19$ as a robust range for
the effect of insurance on blood-sugar screening, with the lower end the more
defensible estimate once positivity is enforced.

\newpage

\section{Acknowledgments}

\textbf{Sujit Silas Armstrong} assembled and cleaned the analytic dataset
(\texttt{nhanes\_data\_prep.Rmd}), implemented the doubly robust (AIPW)
estimator (\texttt{modeling/doubly\_robust.Rmd}), and conducted the sensitivity
analyses (\texttt{sensitivity/sensitivity\_analysis.Rmd}). \textbf{Janvi
Bharucha} wrote the introduction and dataset description and carried out the
survey-weighted regression analysis (\texttt{modeling/outcome\_regression.Rmd}).
\textbf{Amanda Millatt} constructed the DAG and identification strategy
(\texttt{nhanes\_dag.Rmd}) and performed the inverse-probability-weighting
analysis (\texttt{modeling/ipw.Rmd}).

Project repository: \url{https://github.com/sujitsilas/m235-causal-inference-final-project}.

\newpage

\bibliographystyle{plain}  % or "apalike", "ieeetr", "acm", etc.
\bibliography{references}

\end{document}
